\documentclass{article}

\usepackage[utf8]{inputenc}
\usepackage[T1]{fontenc}
\usepackage{helvet}
\usepackage{geometry}
\usepackage{xcolor}
\usepackage{eso-pic}
\usepackage{fancyhdr}
\usepackage{parskip}
\usepackage{microtype}
\usepackage[english, ukrainian]{babel}
\usepackage{url}
\usepackage{amsmath}
\usepackage{amssymb}
\usepackage{amsthm}
\usepackage{longtable}
\usepackage{booktabs}
\usepackage{hyperref}
\usepackage{titlesec}

\definecolor{nato}{HTML}{293757}
\definecolor{footergray}{gray}{0.65}

\geometry{a4paper,left=3cm,right=3cm,top=3cm,bottom=3cm}
\renewcommand{\familydefault}{\sfdefault}
\setcounter{secnumdepth}{3}

\titleformat{\section}
  {\color{nato}\Huge\bfseries}{\thesection.}{1em}{}
\titlespacing*{\section}{0pt}{30pt}{12pt}

\titleformat{\subsection}
  {\color{nato}\LARGE\bfseries}{\thesubsection.}{1em}{}
\titlespacing*{\subsection}{0pt}{20pt}{8pt}

\titleformat{\subsubsection}
  {\color{nato}\Large\bfseries}{\thesubsubsection.}{1em}{}
\titlespacing*{\subsubsection}{0pt}{10pt}{6pt}

\fancyhf{}
\fancyhead[R]{\small\color{footergray} 31 березня 2026}
\fancyfoot[L]{\small\color{footergray} Технічні вимоги ЄСІКС. IV. Реєстри}

\fancyfoot[R]{\small\color{footergray}\thepage}

\newcommand\HUGE[1]{\fontsize{40}{40}\selectfont #1}

\renewcommand{\headrulewidth}{0pt}
\renewcommand{\footrulewidth}{0.4pt}
\renewcommand{\footrule}{\color{footergray}\hrule width \textwidth height 0.4pt}

\begin{document}

\pagestyle{fancy}

\begin{titlepage}
  \AddToShipoutPictureBG*{\AtPageLowerLeft{\color{nato}\rule{\paperwidth}{\paperheight}}}
  \color{white}\thispagestyle{empty}\vspace*{4cm}\centering
  {\Huge  \textbf{Технічні вимоги ЄСІКС} \par} \vspace{2cm}
  {\HUGE  \textbf{IV. Реєстри} \par} \vspace{2.5cm}
  {\LARGE \textbf{IV-1. Виконавчі документи} \par}
  {\LARGE \textbf{IV-2. Судові рішення} \par}
  {\LARGE \textbf{IV-3. Правові позиції} \par}
  {\LARGE \textbf{IV-4. Дайджести} \par}
  {\LARGE \textbf{IV-5. Фізичні особи} \par}
  {\LARGE \textbf{IV-6. Адвокати} \par}
  {\LARGE \textbf{IV-7. Уповноважені} \par}
  {\LARGE \textbf{IV-8. Кадри судової влади} \par}
  {\LARGE \textbf{IV-9. Суб'єкти господарювання} \par}
  {\LARGE \textbf{IV-10. Зовнішні реєстри (ШБО)} \par}
  \vfill
  {\large 2026 \copyright\ Інформаційні Судові Системи \par}
\end{titlepage}

\newpage
\begin{center}
  {\Huge\color{nato}\bfseries Зміст\par}
\end{center}
\vspace{40pt}
\tableofcontents

\begin{abstract}
ТЕХНІЧНІ ВИМОГИ
НА СТВОРЕННЯ ЗАСОБУ ІНФОРМАТИЗАЦІЇ - КОМПОНЕНТІВ “РЕЄСТРИ”, МАЙСТЕР РЕЄСТРІВ, ЄДРСР, ЄДРВД В РАМКАХ
ЄДИНОЇ СУДОВОЇ ІНФОРМАЦІЙНО-КОМУНІКАЦІЙНОЇ СИСТЕМИ.
\end{abstract}

\newpage
\section*{Вступна частина}

\subsection*{Умовні скорочення та визначення}

Терміни та скорочення, що використовуються в документі:

\begin{longtable}{p{4cm}p{10cm}}
\toprule
\textbf{Термін} & \textbf{Значення} \\
\midrule
API & Автоматизовані програмні інтеграції, відомі також як “прикладний інтерфейс програмування – application programming interface” \\
БД & База Даних — Організована колекція даних, яка зберігається та управляється з використанням комп'ютерної системи; забезпечує зберігання, організацію та доступ до даних з метою ефективного використання та обробки \\
БП & Бізнес Процес — Сукупність дій необхідних для здійснення людиною або системою, в результаті яких відбувається досягнення бажаного результату \\
КЕП & Кваліфікований електронний підпис (X.509) \\
ЄСІКС & Єдина судова інформаційно-комунікаційна система \\
\bottomrule
\end{longtable}

\subsection*{Загальні положення}

Технічні вимоги на основі «Концепції ЄСІКС» від 2024 року \cite{esiks} і API судової системи ЄДРСР\cite{eCourt}.

\newpage
\section{Реєстр виконавчих документів}

\subsection{Вимоги до сутностей}

\subsubsection{Виконавчі листи}
\subsubsection{Судові накази}
\subsubsection{Виконавчі написи нотаріусів}
\subsubsection{Ухвали}
\subsubsection{Посвідчення комісій по трудових спорах}
\subsubsection{Постанови органів}
\subsubsection{Постанови державних виконавців}
\subsubsection{Постанови приватних виконавців}
\subsubsection{Рішення і акти виконавчих органів}
\subsubsection{Рішення Європейського суду з прав людини}
\subsubsection{Рішення (постанови) суб’єктів державного фінансового моніторингу}
\subsubsection{Рішення Органу суспільного нагляду за аудиторською діяльністю або Аудиторської палати України}
\subsubsection{Рішення Національної ради України з питань телебачення і радіомовлення}

\subsection{Вимоги до процесів}

\subsubsection{Публікація наказу}
\subsubsection{Публікація ухвали}
\subsubsection{Публікація постанови}
\subsubsection{Публікація рішення}
\subsubsection{Життєвий цикл посвічення}
\subsubsection{Життєвий цикл виконавчого напису}

\newpage
\section{Реєстр судових рішень}
Технічні вимоги до реєстру судових рішень в повній мірі реалізуються ЄДРСР (версія 1.16).

\subsection{Вимоги до довідників}

\subsubsection{Категорія справи (CaseCategory)}
\subsubsection{Роль судді у справі / провадженні (CaseJudgeRole)}
\subsubsection{Тип (роль) учасника справи (CaseMemberRole)}
\subsubsection{Стан розгляду справи (CaseStatus)}
\subsubsection{Суд (Court)}
\subsubsection{Тип суду (CourtType)}
\subsubsection{Категорія документу (DocCategory)}
\subsubsection{Тип документу (DocType)}
\subsubsection{Юрисдикція суду (JurisdictionType)}
\subsubsection{Вид судочинства (JusticeType)}
\subsubsection{Результат розгляду провадження (ProcResult)}
\subsubsection{Стадія розгляду провадження (ProcStage)}
\subsubsection{Стан розгляду провадження (ProcStatus)}
\subsubsection{Регіон (Region)}
\subsubsection{Тип заяви (ClaimType)}

\addtocontents{toc}{\protect\newpage}
\newpage
\subsection{Вимоги до сутностей}

\subsubsection{Судові рішення (DocumentDto)}
\subsubsection{Файли документу (DocumentFileDto)}
\subsubsection{Ідентифікаційні документи (DocumentAsdsLink)}
\subsubsection{Заява (ClaimPostDto)}
\subsubsection{Квинтація заяви (ClaimTicketDto)}
\subsubsection{Справа (Case)}
\subsubsection{Провадження (Proceedings)}
\subsubsection{Учасник справи (CaseMemberDto)}
\subsubsection{Суддя по справі (CaseJudge)}
\subsubsection{Календар судових засідань (Trial)}
\subsubsection{Виконавче провадження (ExecutiveDecision)}
\subsubsection{Довіреність (PowerOfAttorney)}
\subsubsection{Документи сторін (PartyDocPostDto)}
\subsubsection{Постанова виконавчого провадження (ExecutiveDecisionDocumentDto)}
\subsubsection{Результат розсилки документів сторонам (PartyDocPostResultDto)}
\subsubsection{Квитанція про доставку cтороні справи (CaseMemberTicketDto)}

\newpage
\subsection{Вимоги до процесів}

\subsubsection{Завантаження файлу}
\subsubsection{Видалення файлу}
\subsubsection{Життєвий цикл справ}
\subsubsection{Життєвий цикл заяви}
\subsubsection{Життєвий цикл довіреності}
\subsubsection{Життєвий цикл виконавчого провадження}
\subsubsection{Розсилка документів сторонам}
\subsubsection{Отримання квитанції про розсилку сторонам}
\subsubsection{Подача заяви в чергу на обробку}
\subsubsection{Отримання квитанції про статус заяви}
\subsubsection{Публікація судового рішення}
\subsubsection{Перевірка наявності кабінету особи}

\newpage
\section{Реєстр правових позицій}

\subsection{Вимоги до сутностей}

\subsubsection{Правова позиція}

\subsection{Вимоги до процесів}

\subsubsection{Життєвий цикл правової позиції}

\newpage
\section{Реєстр дайджестів}

\subsection{Вимоги до сутностей}

\subsubsection{Дайджест}

\subsection{Вимоги до процесів}

\subsubsection{Життєвий цикл дайджеста}

\addtocontents{toc}{\protect\newpage}
\newpage
\section{Реєстр фізичних осіб}

\subsection{Вимоги до сутностей}

\subsubsection{Фізична особа}

\subsection{Вимоги до процесів}

\subsubsection{Життєвий цикл фізичної особи}

\newpage
\section{Реєстр адвокатів}

\subsection{Вимоги до сутностей}

\subsubsection{Адвокат}

\subsection{Вимоги до процесів}

\subsubsection{Життєвий цикл адвоката}

\newpage
\section{Реєстр суддів}

\subsection{Вимоги до сутностей}

\subsubsection{Суддя}

\subsection{Вимоги до процесів}

\subsubsection{Життєвий цикл судді}

\newpage
\section{Реєстр уповноважених}

\subsection{Вимоги до сутностей}

\subsubsection{Уповноважений}

\subsection{Вимоги до процесів}

\subsubsection{Життєвий цикл уповноваженого}

\newpage
\section{Реєстр суб'єктів Мінекономіки}

\subsection{Вимоги до сутностей}

\subsubsection{Суб'єкт господарської діяльності}

\subsection{Вимоги до процесів}

\subsubsection{Життєвий цикл суб'єкта господарювання}

\newpage
\section*{Висновок}

\begin{thebibliography}{9}

\bibitem{ecoz:rehab} Міністерство охорони здоров’я України. \newblock Технічні вимоги Реабілітація 2.0. \newblock Версія 08.08.2024.
\bibitem{dstu:3973} ДСТУ 3973-2000. \newblock Система розроблення та постачання продукції на виробництво. Правила виконання науково-дослідних робіт. Загальні положення.
\bibitem{dstu:3974} ДСТУ 3974-2000. \newblock Система розроблення та постачання продукції на виробництво. Правила виконання науково-конструкторських робіт. Загальні положення.
\bibitem{dstu:42010} ДСТУ 42010:2018. \newblock Інженерія систем і програмних засобів. Опис архітектури.
\bibitem{law:info-protection} Закон України «Про захист інформації в інформаційно-телекомунікаційних системах».
\bibitem{law:personal-data} Закон України «Про захист персональних даних».
\bibitem{law:cybersecurity} Закон України «Про основні засади забезпечення кібербезпеки України».
\bibitem{owasp:top10} OWASP Foundation. \newblock OWASP Top 10: The Ten Most Critical Web Application Security Risks. \newblock \url{https://owasp.org/www-project-top-ten/}.
\bibitem{nist:sp800-57} NIST Special Publication 800-57 Part 1 Revision 5. \newblock Recommendation for Key Management.
\bibitem{iso:42010} ISO/IEC/IEEE 42010:2022. \newblock Systems and software engineering — Architecture description.
\bibitem{zachman} Zachman, J.A. \newblock A Framework for Information Systems Architecture. \newblock IBM Systems Journal, 1987.
\bibitem{esiks} Концепція ЄСІКС. \url{https://court.gov.ua/storage/portal/dsa/news/Kontseptsia%20ESIKS.pdf}. 2024
\bibitem{cbor} CBOR/COSE encoding/protocol. \url{https://tonpa.guru/stream/2024/2024-11-20%20CBOR%20COSE.htm}. 2025

\end{thebibliography}

\end{document}
